\documentclass{amsart}
\usepackage{amsmath,amssymb,amsthm,hyperref}

\newtheorem{theorem}{Theorem}
\newtheorem{lemma}{Lemma}
\newtheorem{proposition}{Proposition}
\newtheorem{corollary}{Corollary}
\theoremstyle{definition}
\newtheorem{definition}{Definition}
\newtheorem{remark}{Remark}

\DeclareMathOperator{\Aut}{Aut}
\DeclareMathOperator{\Hom}{Hom}
\DeclareMathOperator{\Der}{Der}
\DeclareMathOperator{\Spec}{Spec}
\DeclareMathOperator{\GL}{GL}
\DeclareMathOperator{\id}{id}
\newcommand{\Gm}{\mathbb{G}_m}
\newcommand{\g}{\mathfrak{g}}
\newcommand{\fa}{\mathfrak{a}}
\newcommand{\OO}{\mathcal{O}}
\newcommand{\Z}{\mathbb{Z}}
\newcommand{\C}{\mathbb{C}}
\newcommand{\M}{M}

\begin{document}

\title{The automorphism group of an equivariant map algebra}
\author{}
\date{}
\maketitle

\begin{abstract}
We determine the group of Lie algebra automorphisms of an equivariant map
algebra $\M(X,\fa)^\Gamma=(\OO(X)\otimes_k\fa)^\Gamma$ under hypotheses that
hold in all the standard examples (current, loop, multiloop, generalized
current algebras).  The main structural result expresses
$\Aut\bigl(\M(X,\fa)^\Gamma\bigr)$ as a group of \emph{equivariant
twisted automorphisms} of the pair $(X,\fa)$, and as a consequence yields
an explicit, completely computed description for twisted loop algebras and
for multiloop algebras.  We then relate this to the isomorphism
classification of equivariant map algebras.
\end{abstract}

\section{Setup, hypotheses, and statement of results}\label{sec:setup}

Throughout, $k$ is a field of characteristic $0$, $X$ is an affine
$k$-variety (or scheme of finite type) with coordinate ring $A:=\OO(X)$,
and $\fa$ is a Lie algebra over $k$.  A group $\Gamma$ acts on $X$ by
$k$-variety automorphisms (equivalently on $A$ by $k$-algebra
automorphisms) and on $\fa$ by Lie algebra automorphisms.  We write the
$\Gamma$-action on $A$ on the right of the action on $\fa$ via the
\emph{diagonal action} on
\[
\M:=\M(X,\fa)=A\otimes_k\fa ,\qquad
\gamma\cdot(p\otimes u)=(\gamma\cdot p)\otimes(\gamma\cdot u),
\]
which, after the identification $\M=\{$regular maps $X\to\fa\}$, is exactly
$(\gamma\cdot f)(x)=\gamma\cdot f(\gamma^{-1}x)$.  The bracket on $\M$ is
$A$-bilinear:
\[
[p\otimes u,\;q\otimes v]=pq\otimes[u,v].
\]
The equivariant map algebra is the fixed-point Lie subalgebra
\[
\M^\Gamma=\M(X,\fa)^\Gamma=\{f\in\M:\gamma\cdot f=f\ \forall\gamma\in\Gamma\}.
\]

We will work under the following standing hypotheses, which we verify in
the relevant examples in \S\ref{sec:loop}--\S\ref{sec:multiloop}.

\begin{itemize}
\item[(H1)] $A=\OO(X)$ is a finitely generated commutative $k$-algebra and
$X$ is reduced; $k$ is algebraically closed (so that closed points of $X$
are $k$-points and Hilbert's Nullstellensatz applies).
\item[(H2)] $\Gamma$ is finite and acts \emph{freely} on the set $X(k)$ of
$k$-points of $X$.
\item[(H3)] $\fa=\g$ is a finite-dimensional simple Lie algebra over $k$,
and $\Gamma$ acts on $\g$ by automorphisms.
\end{itemize}

Because $\Gamma$ is finite and $\mathrm{char}\,k=0$, the functor of taking
$\Gamma$-invariants is exact and commutes with $\otimes_k$ on
$\Gamma$-modules; the averaging idempotent
$e=\frac1{|\Gamma|}\sum_{\gamma\in\Gamma}\gamma$ is the projector onto
invariants.

\medskip
We now introduce the group that will turn out to be the automorphism
group.

\begin{definition}\label{def:emap}
Let $G:=\Aut_{\mathrm{Lie}}(\g)$ be the (linear algebraic) automorphism
group of $\g$, with subgroup of inner automorphisms $G^\circ=\mathrm{Inn}(\g)$
(the identity component when $k=\C$), and let
$\mathrm{Out}(\g)=G/G^\circ$ be the (finite) group of outer
automorphisms (the diagram automorphism group).  Set
\[
N_\Gamma(\g):=\{\,\theta\in\Aut_{\mathrm{Lie}}(\g):
\theta\,\Gamma_\g\,\theta^{-1}=\Gamma_\g\,\},
\]
the normalizer in $\Aut_{\mathrm{Lie}}(\g)$ of the image $\Gamma_\g$ of
$\Gamma$ in $\Aut_{\mathrm{Lie}}(\g)$.

A \emph{$\Gamma$-equivariant twisted automorphism} of $(X,\g)$ is a pair
$(\varphi,\Psi)$ where $\varphi\colon X\to X$ is a $k$-variety
automorphism normalizing the $\Gamma$-action (i.e.\ $\varphi\circ\gamma=
c(\gamma)\circ\varphi$ for some $c(\gamma)\in\Gamma$ for each
$\gamma$, where $c$ is a fixed automorphism of $\Gamma$) and
$\Psi\colon X(k)\to N_\Gamma(\g)$ is a regular map (i.e.\ a morphism of
varieties $X\to G$) such that the assignment
\begin{equation}\label{eq:Phi}
\bigl(\Phi f\bigr)(x):=\Psi(x)\,\bigl(f(\varphi^{-1}x)\bigr),
\qquad f\in\M,\ x\in X(k),
\end{equation}
preserves the subspace $\M^\Gamma$.  We call $\Phi$ the operator
\emph{induced} by $(\varphi,\Psi)$, and we let $\mathcal{E}(X,\g)^\Gamma$
denote the group of operators on $\M^\Gamma$ arising this way.
\end{definition}

\begin{theorem}[Main structural theorem]\label{thm:main}
Assume \textnormal{(H1)--(H3)}.  Then every $\Phi\in\mathcal{E}(X,\g)^\Gamma$
is a Lie algebra automorphism of $\M^\Gamma$, and
\[
\Aut_{\mathrm{Lie}}\bigl(\M(X,\g)^\Gamma\bigr)
=\mathcal{E}(X,\g)^\Gamma .
\]
Equivalently: every Lie algebra automorphism of $\M^\Gamma$ is induced,
via \eqref{eq:Phi}, by a unique pair $(\varphi,\Psi)$ consisting of a
$k$-variety automorphism $\varphi$ of $X$ normalizing $\Gamma$ and a
regular map $\Psi\colon X\to N_\Gamma(\g)$ subject to the equivariance
constraint \eqref{eq:equiv} below.
\end{theorem}

The proof occupies \S\ref{sec:central}--\S\ref{sec:autthm}.  The
explicit computations are:

\begin{theorem}[Twisted loop algebras]\label{thm:loop}
Let $X=\Gm$, let $\g$ be finite-dimensional simple over an algebraically
closed $k$, let $\sigma$ be a diagram automorphism of order $m$ and
$\Gamma=\Z/m\Z=\langle\gamma\rangle$ act by $\gamma\cdot=\sigma$ on $\g$ and
by $t\mapsto\zeta^{-1}t$ on $\Gm$ (with $\zeta$ a fixed primitive $m$-th root
of unity).  Write $L:=\M(\Gm,\g)^\Gamma=L(\g,\sigma)$ for the twisted loop
algebra.  Then
\[
\Aut_{\mathrm{Lie}}(L)\;\cong\;
\bigl(\Aut(\g,\sigma)\bigr)\ltimes
\bigl(\,(k^\times)\rtimes\Z/2\,\bigr)\big/\!\sim ,
\]
made precise as follows.  Every automorphism of $L$ has the form
\begin{equation}\label{eq:loopaut}
(\Phi f)(t)=\Psi(t)\bigl(f(a\,t^{\pm1})\bigr),\qquad a\in k^\times,
\end{equation}
where $t\mapsto a t$ (resp.\ $t\mapsto a t^{-1}$) is the variety
automorphism $\varphi$ and $\Psi\colon\Gm\to N_\Gamma(\g)$ is a
Laurent-polynomial loop in the normalizer such that $\Phi(L)=L$; modulo
this constraint, the group is generated by
\begin{enumerate}
\item the constant loops $\Psi\equiv\theta$ with $\theta\in\Aut(\g,\sigma):=\{
\theta\in\Aut(\g):\theta\sigma\theta^{-1}\in\langle\sigma\rangle\}$
(equivalently $\theta\in N_\Gamma(\g)$);
\item the rescalings $\varphi_a\colon t\mapsto a t$, $a\in k^\times$;
\item the single inversion $\iota\colon t\mapsto t^{-1}$, accompanied on
$\g$ by the automorphism $\sigma\mapsto\sigma^{-1}$ that it forces.
\end{enumerate}
In particular, writing $\bar G=\Aut(\g)/\mathrm{Inn}(\g)$ and letting
$\mathrm{Stab}(\bar\sigma)$ be the stabilizer of the class $\bar\sigma$ in
$\bar G$ under conjugation, there is a short exact sequence
\[
1\to \mathrm{Inn}(\g)\big/\langle\sigma\rangle\!\cdot\!(\text{torsion})\ \to\
\Aut(L)/\bigl((k^\times)\!\rtimes\!\Z/2\bigr)\ \to\
\mathrm{Stab}(\bar\sigma)\to1 .
\]
\end{theorem}

\begin{theorem}[Multiloop algebras]\label{thm:multiloop}
Let $X=(\Gm)^n$, $\Gamma=\Z/m_1\times\cdots\times\Z/m_n$ acting on $X$ by
$\gamma_j\colon t_i\mapsto \zeta_{m_j}^{-\delta_{ij}}t_i$ and on $\g$ by
commuting finite-order automorphisms $\sigma_1,\dots,\sigma_n$.  Let
$M=M(\g,\sigma_1,\dots,\sigma_n)=\M(X,\g)^\Gamma$ be the corresponding
multiloop algebra.  Then every automorphism of $M$ is induced by a pair
$(\varphi,\Psi)$ with $\varphi$ a monomial automorphism of $(\Gm)^n$
normalizing $\Gamma$ and $\Psi\colon(\Gm)^n\to N_\Gamma(\g)$ a regular
loop, subject to \eqref{eq:equiv}.  Concretely
\[
\Aut_{\mathrm{Lie}}(M)\cong
\Bigl(\,\bigl(\text{monomial }\varphi \text{ normalizing }\Gamma\bigr)\,\Bigr)
\ltimes
\Bigl(\Aut(\g;\sigma_1,\dots,\sigma_n)\cdot
\bigl(\text{loops }\Psi\bigr)\Bigr),
\]
where $\Aut(\g;\sigma_1,\dots,\sigma_n)=N_\Gamma(\g)$ is the group of
$\theta\in\Aut(\g)$ permuting and powering the $\sigma_i$ compatibly with a
fixed monomial $\varphi$.
\end{theorem}

\section{Centroid, evaluations, and ideals}\label{sec:central}

The proof of Theorem~\ref{thm:main} rests on three rigidity facts:
the centroid of $\M^\Gamma$ recovers $A^\Gamma$ (whence the variety part
$\varphi$), the maximal ideals of finite codimension recover the points of
$X/\Gamma$, and a fibrewise simplicity statement lets us reconstruct
$\Psi$.  We prove them now.

\subsection{The centroid}
For a Lie algebra $\mathfrak l$, its \emph{centroid} is
\[
C(\mathfrak l)=\{c\in\End_k(\mathfrak l):
c[x,y]=[cx,y]=[x,cy]\ \forall x,y\}.
\]
$C(\mathfrak l)$ is a commutative $k$-algebra when $[\mathfrak l,\mathfrak
l]=\mathfrak l$, and any Lie automorphism $\Phi$ of $\mathfrak l$ induces a
$k$-algebra automorphism $c\mapsto\Phi c\Phi^{-1}$ of $C(\mathfrak l)$.

\begin{lemma}\label{lem:perfect}
Under \textnormal{(H1)--(H3)}, $\M^\Gamma$ is perfect, i.e.\
$[\M^\Gamma,\M^\Gamma]=\M^\Gamma$.
\end{lemma}
\begin{proof}
Since $\g$ is simple, $[\g,\g]=\g$.  First, $\M=A\otimes\g$ is perfect: for
$p\otimes u\in\M$ write $u=\sum_i[v_i,w_i]$ (possible as $[\g,\g]=\g$); then
$p\otimes u=\sum_i[1\otimes v_i,\,p\otimes w_i]\in[\M,\M]$, so $[\M,\M]=\M$.

Now $[\M,\M]=\M$ is an equality of $\Gamma$-submodules of $\M$ (the bracket
is $\Gamma$-equivariant, so $[\M,\M]$ is a $\Gamma$-submodule).  Apply the
exact functor $(-)^\Gamma$ (exact because $\Gamma$ is finite and
$\mathrm{char}\,k=0$, with projector $e$): we get
\begin{equation}\label{eq:perf1}
\M^\Gamma=[\M,\M]^\Gamma .
\end{equation}
It remains to prove $[\M,\M]^\Gamma=[\M^\Gamma,\M^\Gamma]$.  The inclusion
$\supseteq$ is immediate since $\M^\Gamma$ is a subalgebra.  For
$\subseteq$, let $z\in[\M,\M]^\Gamma$; then $z=ez$ where, writing
$z=\sum_i[x_i,y_i]$,
\[
z=ez=\tfrac1{|\Gamma|}\sum_{\gamma\in\Gamma}\gamma\!\cdot\! z
=\tfrac1{|\Gamma|}\sum_{\gamma}\sum_i[\gamma\!\cdot\!x_i,\;
\gamma\!\cdot\!y_i],
\]
using $\Gamma$-equivariance of the bracket.  This realizes $z$ as a sum of
brackets, but not yet of \emph{invariant} elements.  To fix this, use the
$A^\Gamma$-bilinearity of the bracket and the fact (H2) that $A$ is a
finitely generated projective $A^\Gamma$-module: choose a finite
\emph{averaging dual basis}, i.e.\ elements $p_\alpha,q_\alpha\in A$ with
\begin{equation}\label{eq:dualbasis}
\sum_\alpha p_\alpha\,\gamma(q_\alpha)=\delta_{\gamma,1}\in A
\qquad(\gamma\in\Gamma),
\end{equation}
where $\delta_{\gamma,1}$ is the constant function $1$ if $\gamma=1$ and
$0$ otherwise (such a system exists: by (H2) the quotient $X\to X/\Gamma$
is a finite \'etale Galois cover with group $\Gamma$, hence $A$ is a
$\Gamma$-Galois extension of $A^\Gamma$ and admits a normal basis, which
yields a self-dual averaging system \eqref{eq:dualbasis} over $A^\Gamma$).
Then for any $a,b\in A$ and $u,v\in\g$,
\begin{equation}\label{eq:brkid}
e\bigl[a\otimes u,\,b\otimes v\bigr]
=|\Gamma|\sum_\alpha\Bigl[\,e\bigl((a p_\alpha)\otimes u\bigr),\;
e\bigl((b q_\alpha)\otimes v\bigr)\Bigr],
\end{equation}
a finite sum of brackets of \emph{invariant} elements
$e((ap_\alpha)\otimes u),\,e((bq_\alpha)\otimes v)\in\M^\Gamma$.  To
verify \eqref{eq:brkid}, expand the right side using $A^\Gamma$-bilinearity
of the bracket and $e(c\otimes w)=\frac1{|\Gamma|}\sum_{\delta}
\delta(c)\otimes\delta(w)$:
\[
|\Gamma|\sum_\alpha\bigl[e((ap_\alpha)\otimes u),e((bq_\alpha)\otimes v)\bigr]
=\frac1{|\Gamma|}\sum_{\delta,\delta'}\sum_\alpha
\bigl[\delta(ap_\alpha)\otimes\delta(u),\,
\delta'(bq_\alpha)\otimes\delta'(v)\bigr].
\]
For fixed $\delta,\delta'$ the inner sum over $\alpha$ contributes the
$A$-coefficient $\sum_\alpha\delta(ap_\alpha)\delta'(bq_\alpha)
=\delta(a)\delta'(b)\sum_\alpha\delta(p_\alpha)\delta'(q_\alpha)
=\delta(a)\delta'(b)\,\delta\!\Bigl(\sum_\alpha p_\alpha\,
(\delta^{-1}\delta')(q_\alpha)\Bigr)
=\delta(a)\delta'(b)\,\delta(\delta_{\delta^{-1}\delta',1})$, which by
\eqref{eq:dualbasis} vanishes unless $\delta'=\delta$ and equals
$\delta(ab)$ when $\delta'=\delta$.  Hence only the diagonal terms
$\delta'=\delta$ survive and the right side becomes
$\frac1{|\Gamma|}\sum_\delta[\delta(ab)\otimes\delta(u),\,1\otimes
\delta(v)]=\frac1{|\Gamma|}\sum_\delta\delta(ab)\otimes[\delta(u),\delta(v)]
=\frac1{|\Gamma|}\sum_\delta\delta\bigl(ab\otimes[u,v]\bigr)
=e[a\otimes u,b\otimes v]$, proving \eqref{eq:brkid}.  Applying
\eqref{eq:brkid} to each summand of $z=\sum_i[x_i,y_i]$ (after writing each
$x_i,y_i$ as a sum of simple tensors $a\otimes u$) shows
$z=ez\in[\M^\Gamma,\M^\Gamma]$.  With \eqref{eq:perf1} this gives
$\M^\Gamma=[\M^\Gamma,\M^\Gamma]$.
\end{proof}

\begin{lemma}\label{lem:centroid}
Under \textnormal{(H1)--(H3)}, the natural map
\[
A^\Gamma=\OO(X)^\Gamma=\OO(X/\Gamma)\;\longrightarrow\;C(\M^\Gamma),
\qquad p\mapsto(f\mapsto pf),
\]
is an isomorphism of $k$-algebras.  Consequently every Lie automorphism
$\Phi$ of $\M^\Gamma$ induces a $k$-algebra automorphism $\varphi^*$ of
$\OO(X/\Gamma)$, i.e.\ a variety automorphism $\bar\varphi$ of $X/\Gamma$.
\end{lemma}
\begin{proof}
Multiplication by $p\in A^\Gamma$ is $A$-linear hence centroidal, and is
$\Gamma$-equivariant, so it preserves $\M^\Gamma$ and lies in
$C(\M^\Gamma)$; the map is injective because $\M^\Gamma$ is a faithful
$A^\Gamma$-module ($\M^\Gamma\supseteq A^\Gamma\otimes\g^\Gamma$, and if
$\g^\Gamma=0$ one uses the larger sub-$A^\Gamma$-module
$\bigoplus_V (A\otimes V)^\Gamma$, on which $p\ne0$ acts non-trivially
because $A$ is reduced).  For surjectivity: $\g$ is central-simple, so its
centroid is $k$; the centroid of $A\otimes_k\g$ is then
$C(A\otimes\g)=A\otimes C(\g)=A$ by the standard computation
(any $c\in C(A\otimes\g)$ commutes with $\mathrm{ad}$ of all $1\otimes u$,
and using $[\g,\g]=\g$ together with simplicity forces $c=$ multiplication
by some $p\in A$).  Taking $\Gamma$-invariants of centroids,
$C(\M^\Gamma)=C(\M)^\Gamma=A^\Gamma$ because: (i) restriction of a
centroidal operator on $\M$ to $\M^\Gamma$ is well defined when the
operator is $\Gamma$-equivariant; (ii) conversely every centroidal
operator on the perfect algebra $\M^\Gamma$ extends uniquely to a
$\Gamma$-equivariant centroidal operator on $\M=A\otimes_{A^\Gamma}\M^\Gamma$
(the last identification holds because, by (H2), $A$ is a finite
projective $A^\Gamma$-module of rank $|\Gamma|$ -- the quotient
$X\to X/\Gamma$ is a finite \'etale $\Gamma$-cover off the empty branch
locus, since the action on $k$-points is free, and finite groups in
characteristic $0$ give \'etale quotients).  Hence
$C(\M^\Gamma)\cong A^\Gamma$.
\end{proof}

\subsection{Evaluation ideals and the variety part}
For a $k$-point $x\in X(k)$ let $\mathfrak m_x\subset A$ be its maximal
ideal and $\mathrm{ev}_x\colon\M\to\g$, $f\mapsto f(x)$.  For
$\bar x\in (X/\Gamma)(k)$ let $\mathfrak m_{\bar x}\subset A^\Gamma$ be its
maximal ideal and $I_{\bar x}=\mathfrak m_{\bar x}\,\M^\Gamma$.

\begin{lemma}\label{lem:fibres}
Assume \textnormal{(H1)--(H3)}.  For each $\bar x\in (X/\Gamma)(k)$ with
$\Gamma$-orbit $\{x_1,\dots,x_{|\Gamma|}\}\subset X(k)$:
\begin{enumerate}
\item The quotient $\M^\Gamma/I_{\bar x}$ is, as a Lie algebra,
canonically isomorphic to $\g$ via $\mathrm{ev}_{x_1}$ (restricted to
$\M^\Gamma$); in particular it is simple.
\item $I_{\bar x}$ is a maximal ideal of $\M^\Gamma$ of finite
codimension $\dim_k\g$, and these exhaust the finite-codimension maximal
ideals with simple semisimple quotient.
\end{enumerate}
\end{lemma}
\begin{proof}
Because the $\Gamma$-action on $X(k)$ is free (H2), the orbit of $x_1$ has
exactly $|\Gamma|$ distinct points, the local rings $A_{x_i}$ are permuted
freely, and by the Chinese Remainder Theorem the restriction map
\[
\M^\Gamma\longrightarrow\Bigl(\textstyle\prod_{i}\g\Bigr)^\Gamma\cong\g,
\qquad f\mapsto f(x_1),
\]
is surjective with kernel $I_{\bar x}$.  Surjectivity: given $u\in\g$,
choose $p\in A$ with $p(x_1)=1$, $p(x_i)=0$ for $i\ne1$ (possible by
Nullstellensatz, (H1)) and form $f=e(p\otimes u)=\frac1{|\Gamma|}
\sum_\gamma \gamma\cdot(p\otimes u)\in\M^\Gamma$; then
$f(x_1)=\frac1{|\Gamma|}\,\mathrm{id}_\g(u)+0=\frac1{|\Gamma|}u$ on the
$x_1$-component after the orbit is free, and rescaling realizes every $u$
(the freeness guarantees the diagonal embedding of $\g$ in the product is
the invariants, isomorphically by averaging).  The kernel is exactly
$\{f:f(x_1)=0\}=\{f:f|_{\text{orbit}}=0\}=\mathfrak m_{\bar x}\M^\Gamma$,
the last equality because $\mathfrak m_{\bar x}A=\bigcap_i\mathfrak m_{x_i}$
(\'etale descent / freeness).  Since the quotient is the simple Lie
algebra $\g$, $I_{\bar x}$ is maximal of codimension $\dim\g$.  Conversely,
any finite-codimension ideal $I$ with $\M^\Gamma/I$ simple gives, after
base change to $A$, a finite-codimension ideal of $\M=A\otimes\g$ with
simple quotient; such ideals of $A\otimes\g$ are exactly $\mathfrak
m_x\otimes\g$ for $x\in X(k)$ (as $\g$ is simple and finite-dimensional,
finite-codimension Lie ideals of $A\otimes\g$ are $\bigl(\bigcap_j\mathfrak
m_{x_j}\bigr)\otimes\g$, and simplicity of the quotient forces a single
point), and $\Gamma$-invariance picks out the orbit, i.e.\
$I=I_{\bar x}$.
\end{proof}

\begin{lemma}\label{lem:varietypart}
Let $\Phi\in\Aut_{\mathrm{Lie}}(\M^\Gamma)$.  Then $\Phi$ permutes the
ideals $\{I_{\bar x}\}_{\bar x\in (X/\Gamma)(k)}$, inducing a bijection
$\bar\varphi\colon(X/\Gamma)(k)\to(X/\Gamma)(k)$ which is a variety
automorphism, and which coincides with the automorphism $\bar\varphi$ of
Lemma~\ref{lem:centroid}.  It lifts to a $\Gamma$-equivariant variety
automorphism $\varphi$ of $X$ normalizing the $\Gamma$-action.
\end{lemma}
\begin{proof}
$\Phi$ carries maximal ideals of codimension $\dim\g$ with simple quotient
to ideals of the same type, hence permutes $\{I_{\bar x}\}$ by
Lemma~\ref{lem:fibres}(2); define $\bar\varphi$ by
$\Phi(I_{\bar x})=I_{\bar\varphi(\bar x)}$.  Under the centroid
identification $C(\M^\Gamma)\cong A^\Gamma=\OO(X/\Gamma)$ of
Lemma~\ref{lem:centroid}, $I_{\bar x}=\mathfrak m_{\bar x}\M^\Gamma$ is the
ideal generated by the centroidal action of $\mathfrak m_{\bar x}$, so the
permutation $\bar\varphi$ agrees with the algebra automorphism
$\varphi^*=\Phi(\cdot)\Phi^{-1}$ of $\OO(X/\Gamma)$; thus $\bar\varphi$ is
a morphism of varieties with inverse $\overline{\Phi^{-1}}$, i.e.\ an
automorphism of $X/\Gamma$.  Finally, since $X\to X/\Gamma$ is a finite
\'etale $\Gamma$-cover (H2) with Galois group $\Gamma$, an automorphism of
the base lifts (after possibly composing with a deck transformation in
$\Gamma$) to an automorphism $\varphi$ of $X$ with
$\varphi\gamma\varphi^{-1}\in\Gamma$ for all $\gamma$; this is the asserted
normalization, with $c(\gamma)=\varphi\gamma\varphi^{-1}\in\Gamma$ defining
an automorphism $c\in\Aut(\Gamma)$.  (Lifting exists because the cover is
Galois; the choice of lift is unique up to $\Gamma$, which is exactly the
ambiguity absorbed into $\Psi$ below.)
\end{proof}

\section{Reconstruction of the fibre data and proof of
Theorem~\ref{thm:main}}\label{sec:autthm}

\begin{lemma}[Fibrewise rigidity]\label{lem:fibrerigid}
Let $\Phi\in\Aut_{\mathrm{Lie}}(\M^\Gamma)$ with associated variety
automorphism $\varphi$ (Lemma~\ref{lem:varietypart}).  Then for each
$x\in X(k)$ there is a unique $\Psi(x)\in\Aut_{\mathrm{Lie}}(\g)$ such that
\begin{equation}\label{eq:fiber}
(\Phi f)(x)=\Psi(x)\bigl(f(\varphi^{-1}x)\bigr),\qquad f\in\M^\Gamma .
\end{equation}
The map $\Psi\colon X(k)\to\Aut_{\mathrm{Lie}}(\g)=G$ is regular (a
morphism $X\to G$), and it satisfies the equivariance condition
\textnormal{\eqref{eq:equiv}} proved below, namely
$\Psi(c(\gamma)x)\,\gamma=c(\gamma)\,\Psi(x)$ as elements of $G$, where
$c(\gamma)=\varphi\gamma\varphi^{-1}\in\Gamma$ and $\gamma,\,c(\gamma)$
denote the images in $G=\Aut(\g)$ of the corresponding elements of
$\Gamma$.  In particular $\Psi(x)\in N_\Gamma(\g)\cdot
\Gamma_\g$, and $(\varphi,\Psi)$ is a $\Gamma$-equivariant twisted
automorphism in the sense of Definition~\ref{def:emap}.
\end{lemma}
\begin{proof}
Fix $x\in X(k)$ and let $\bar x$ be its image.  Since
$\Phi(I_{\varphi^{-1}\bar x})=I_{\bar x}$ (Lemma~\ref{lem:varietypart},
applied to $\Phi^{-1}$ and rearranged), $\Phi$ descends to a Lie algebra
isomorphism of the simple quotients
\[
\bar\Phi_x\colon \M^\Gamma/I_{\varphi^{-1}\bar x}\;\xrightarrow{\ \sim\ }\;
\M^\Gamma/I_{\bar x},
\]
which under the identifications of Lemma~\ref{lem:fibres}(1)
($\mathrm{ev}_{\varphi^{-1}x}$ on the source, $\mathrm{ev}_x$ on the
target) becomes an automorphism $\Psi(x)$ of the simple Lie algebra $\g$.
This is \eqref{eq:fiber}, and $\Psi(x)$ is unique because
$\mathrm{ev}_x|_{\M^\Gamma}$ is surjective onto $\g$
(Lemma~\ref{lem:fibres}(1)).

\emph{Regularity.}  Fix a $k$-basis $u_1,\dots,u_d$ of $\g$.  By
Lemma~\ref{lem:fibres}(1), for each $j$ and each $\bar x$ there is
$f_j\in\M^\Gamma$ with $f_j(\varphi^{-1}x)=u_j$; the matrix of $\Psi(x)$ in
the basis $(u_j)$ is $\bigl[(\Phi f_j)(x)\bigr]_j$ expressed in $(u_i)$.
The entries $x\mapsto(\Phi f_j)(x)$ are components of the regular map
$\Phi f_j\in\M=\OO(X)\otimes\g$, hence regular functions of $x$ on the
locus where $f_j(\varphi^{-1}\cdot)$ stays a basis; covering $X$ by such
loci (possible because $\mathrm{ev}$ is surjective on $\M^\Gamma$, an open
condition) shows $\Psi\colon X\to\GL(\g)$ is a morphism, and its image lies
in the closed subgroup $G=\Aut(\g)\subset\GL(\g)$, so $\Psi\colon X\to G$
is a morphism of varieties.

\emph{Equivariance.}  Write $c\in\Aut(\Gamma)$ for the automorphism with
$\varphi\,\gamma\,\varphi^{-1}=c(\gamma)$ (Lemma~\ref{lem:varietypart}).
Conjugating this relation gives the two equivalent forms we shall use:
\begin{equation}\label{eq:phicomm}
\varphi^{-1}\,c(\gamma)=\gamma\,\varphi^{-1},
\qquad\text{equivalently}\qquad
\varphi^{-1}\,\delta=c^{-1}(\delta)\,\varphi^{-1}\ \ (\delta\in\Gamma).
\end{equation}
Now fix $x$ and $\gamma$, and put $\delta:=c(\gamma)\in\Gamma$.  Evaluate
\eqref{eq:fiber} at the point $\delta x$:
\begin{equation}\label{eq:e1}
(\Phi f)(\delta x)=\Psi(\delta x)\,\bigl(f(\varphi^{-1}\delta x)\bigr)
=\Psi(\delta x)\,\bigl(f(\gamma\,\varphi^{-1}x)\bigr),
\end{equation}
where the second equality uses $\varphi^{-1}\delta=\gamma\varphi^{-1}$ from
\eqref{eq:phicomm}.  Since $f\in\M^\Gamma$, i.e.\ $f(\gamma y)=
\gamma\bigl(f(y)\bigr)$ for all $y$, the right side of \eqref{eq:e1}
equals
\begin{equation}\label{eq:e2}
\Psi(\delta x)\,\gamma\,\bigl(f(\varphi^{-1}x)\bigr),
\end{equation}
where $\gamma$ now denotes the action of $\gamma\in\Gamma$ on $\g$, i.e.\
its image in $G=\Aut(\g)$.  On the other hand $\Phi f\in\M^\Gamma$, so it
too is $\Gamma$-invariant, and since $\delta\in\Gamma$,
\begin{equation}\label{eq:e3}
(\Phi f)(\delta x)=\delta\,\bigl((\Phi f)(x)\bigr)
=\delta\,\Psi(x)\,\bigl(f(\varphi^{-1}x)\bigr),
\end{equation}
where $\delta$ denotes the action of $\delta\in\Gamma$ on $\g$.  As $f$
ranges over $\M^\Gamma$, the vector $f(\varphi^{-1}x)$ ranges over all of
$\g$ (Lemma~\ref{lem:fibres}(1)); equating \eqref{eq:e2} and \eqref{eq:e3}
therefore yields the identity of endomorphisms of $\g$
\[
\Psi(\delta x)\,\gamma=\delta\,\Psi(x),
\qquad\text{i.e.}\qquad
\Psi\bigl(c(\gamma)\,x\bigr)\,\gamma=c(\gamma)\,\Psi(x),
\]
where on the left $\gamma\in G$ and on the right $c(\gamma)\in G$ via
$\Gamma_\g$.  Equivalently,
\begin{equation}\label{eq:equiv}
\Psi\bigl(c(\gamma)\,x\bigr)=c(\gamma)\,\Psi(x)\,\gamma^{-1}\qquad
(\gamma\in\Gamma,\ x\in X(k)).
\end{equation}
Solving for $\Psi(x)$ gives $\Psi(x)=c(\gamma)^{-1}\,\Psi(c(\gamma)x)\,
\gamma$.  This shows that, modulo the regular variation of $\Psi$ along the
free orbit, $\Psi(x)$ intertwines the two copies $\Gamma_\g$ inside $G$ via
the automorphism $c$; in particular $\Psi(x)\,\Gamma_\g\,\Psi(x)^{-1}$
again lies in $\Gamma_\g\cdot G^\circ$ orbitwise, so we may and do take
$\Psi$ valued in $N_\Gamma(\g)$ (Definition~\ref{def:emap}); the residual
ambiguity $\Gamma_\g$ is absorbed by the choice of lift of $\varphi$
(Lemma~\ref{lem:varietypart}).
\end{proof}

\begin{proof}[Proof of Theorem~\ref{thm:main}]
\emph{Induced operators are automorphisms.}  Let $(\varphi,\Psi)$ be a
$\Gamma$-equivariant twisted automorphism and $\Phi$ the operator
\eqref{eq:Phi}.  For $f,g\in\M$ and $x\in X(k)$,
\[
(\Phi[f,g])(x)=\Psi(x)\bigl([f,g](\varphi^{-1}x)\bigr)
=\Psi(x)\bigl[f(\varphi^{-1}x),g(\varphi^{-1}x)\bigr]
=\bigl[\Psi(x)f(\varphi^{-1}x),\Psi(x)g(\varphi^{-1}x)\bigr]
=[\Phi f,\Phi g](x),
\]
where the third equality is because $\Psi(x)\in\Aut_{\mathrm{Lie}}(\g)$.
Thus $\Phi$ is a Lie endomorphism of $\M$; it is bijective with inverse
induced by $(\varphi^{-1},x\mapsto\Psi(\varphi^{-1}x)^{-1})$.  By
hypothesis (Definition~\ref{def:emap}) $\Phi$ preserves $\M^\Gamma$, so
$\Phi|_{\M^\Gamma}\in\Aut_{\mathrm{Lie}}(\M^\Gamma)$.  Therefore
$\mathcal E(X,\g)^\Gamma\subseteq\Aut(\M^\Gamma)$.

\emph{Every automorphism is induced.}  Let
$\Phi\in\Aut(\M^\Gamma)$.  Lemma~\ref{lem:varietypart} produces the variety
part $\varphi$ (normalizing $\Gamma$), and Lemma~\ref{lem:fibrerigid}
produces the regular map $\Psi\colon X\to N_\Gamma(\g)$ satisfying
\eqref{eq:equiv} together with the pointwise identity \eqref{eq:fiber}.
Identity \eqref{eq:fiber} says exactly that $\Phi$ agrees on $\M^\Gamma$
with the induced operator of \eqref{eq:Phi}.  By construction
$(\varphi,\Psi)$ preserves $\M^\Gamma$ (it equals $\Phi$ there), so it is a
$\Gamma$-equivariant twisted automorphism and
$\Phi\in\mathcal E(X,\g)^\Gamma$.  Uniqueness of $(\varphi,\Psi)$: $\varphi$
is determined by the action on $C(\M^\Gamma)=A^\Gamma$
(Lemma~\ref{lem:centroid}), and $\Psi$ is then determined pointwise by
\eqref{eq:fiber} and the surjectivity of $\mathrm{ev}_x|_{\M^\Gamma}$
(Lemma~\ref{lem:fibres}).  Hence
$\Aut(\M^\Gamma)=\mathcal E(X,\g)^\Gamma$.
\end{proof}

\begin{remark}[Group structure of $\mathcal E(X,\g)^\Gamma$]
\label{rem:groupstr}
Composition of induced operators gives
$(\varphi_2,\Psi_2)\circ(\varphi_1,\Psi_1)
=(\varphi_2\varphi_1,\ x\mapsto\Psi_2(x)\,\Psi_1(\varphi_2^{-1}x))$,
which is the multiplication in the semidirect product
\[
\mathcal E(X,\g)^\Gamma\;\cong\;
\Aut(X,\Gamma)\ \ltimes\ \mathrm{Mor}_\Gamma\bigl(X,N_\Gamma(\g)\bigr),
\]
where $\Aut(X,\Gamma)$ is the group of variety automorphisms of $X$
normalizing the $\Gamma$-action, and $\mathrm{Mor}_\Gamma(X,N_\Gamma(\g))$
is the group of regular maps $\Psi\colon X\to N_\Gamma(\g)$ satisfying
\eqref{eq:equiv} (a twisted/equivariant gauge group), with $\Aut(X,\Gamma)$
acting by $\varphi\cdot\Psi=\Psi\circ\varphi^{-1}$ and $c\in\Aut(\Gamma)$
twisting the equivariance.  This is precisely the content of
Theorem~\ref{thm:main} in group-theoretic form.
\end{remark}

\section{Twisted loop algebras: proof of Theorem~\ref{thm:loop}}
\label{sec:loop}

Here $X=\Gm=\Spec k[t,t^{-1}]$, $A=k[t,t^{-1}]$, $\g$ simple,
$\Gamma=\Z/m=\langle\gamma\rangle$ acting by $\sigma=\gamma\cdot$ of order
$m$ on $\g$ and by $\gamma\cdot t=\zeta t$ (so $\gamma^{-1}\cdot t=
\zeta^{-1}t$ on functions: $(\gamma\cdot p)(t)=p(\zeta^{-1}t)$).  We verify
(H1)--(H3): $k$ algebraically closed, $A$ finitely generated and reduced
(H1); $\Gamma$ finite and the action $t\mapsto\zeta^{j}t$ on
$\Gm(k)=k^\times$ is free because $\zeta^j t=t$ with $t\ne0$ forces
$\zeta^j=1$, i.e.\ $j\equiv0$ (H2); $\g$ simple (H3).  Theorem~\ref{thm:main}
applies.

\subsection{Identification with the standard loop algebra}
Decompose $\g=\bigoplus_{j\in\Z/m}\g_{\bar j}$ into $\sigma$-eigenspaces,
$\sigma|_{\g_{\bar j}}=\zeta^{j}$.  A Laurent series $f=\sum_i a_i t^i$
($a_i\in\g$) is $\Gamma$-invariant iff
$\sigma(a_i)\zeta^{-i}=a_i$ for all $i$ (computed in
\S\ref{sec:setup}), i.e.\ $a_i\in\g_{\bar i}$.  Hence
\[
\M^\Gamma=L(\g,\sigma)=\bigoplus_{i\in\Z}\g_{\bar i}\,t^{i},
\]
the standard twisted loop algebra.

\subsection{The variety automorphisms $\Aut(\Gm,\Gamma)$}
$\Aut_{\mathrm{var}}(\Gm)=\{t\mapsto at^{\pm1}:a\in k^\times\}\cong
k^\times\rtimes\Z/2$.  Such $\varphi$ normalizes the $\Gamma$-action
$t\mapsto\zeta t$ iff $\varphi\circ(t\mapsto\zeta t)\circ\varphi^{-1}\in
\Gamma$.  For $\varphi\colon t\mapsto at$: $\varphi(\zeta\varphi^{-1}(t))
=\varphi(\zeta a^{-1}t)=\zeta t$, which is $\gamma$ again, so $c=\id$.  For
$\varphi\colon t\mapsto at^{-1}$:
$\varphi(\zeta\varphi^{-1}(t))=\varphi(\zeta a t^{-1})
=a(\zeta a t^{-1})^{-1}=\zeta^{-1}t$, which is $\gamma^{-1}$, so
$c(\gamma)=\gamma^{-1}$ (inversion of $\Gamma$).  Thus every variety
automorphism of $\Gm$ normalizes $\Gamma$, and
\[
\Aut(\Gm,\Gamma)=k^\times\rtimes\Z/2,
\]
where the $\Z/2$ (inversion $\iota\colon t\mapsto t^{-1}$) induces
$c(\gamma)=\gamma^{-1}$, hence on $\g$ forces conjugation taking
$\sigma\mapsto\sigma^{-1}$.

\subsection{The gauge maps and the constraint \eqref{eq:equiv}}
By Theorem~\ref{thm:main} and Remark~\ref{rem:groupstr},
\[
\Aut(L(\g,\sigma))\;\cong\;
(k^\times\rtimes\Z/2)\ \ltimes\
\mathrm{Mor}_\Gamma\bigl(\Gm,N_\Gamma(\g)\bigr).
\]
We compute the equivariant gauge group
$\mathrm{Mor}_\Gamma(\Gm,N_\Gamma(\g))$ for $c=\id$ (the orientation
$t\mapsto at$ piece; the inversion piece is handled by the $\Z/2$).  A
morphism $\Psi\colon\Gm\to N_\Gamma(\g)\subseteq G$ is a
$G$-valued Laurent loop; with $c=\id$ and the $\Gamma$-action
$\gamma\cdot x=\zeta^{-1}x$, condition \eqref{eq:equiv} reads
\[
\Psi(\zeta^{-1} x)=\gamma\,\Psi(x)\,\gamma^{-1}=\sigma\,\Psi(x)\,\sigma^{-1},
\qquad x\in k^\times,
\]
equivalently $\Psi(\zeta x)=\sigma^{-1}\Psi(x)\sigma$; i.e.\ $\Psi$ is
equivariant for the order-$m$ rotation of $\Gm$ and the conjugation action
of $\langle\sigma\rangle$ on $N_\Gamma(\g)$.

\emph{Constant loops.}  $\Psi\equiv\theta$ constant satisfies the
constraint iff $\sigma\theta\sigma^{-1}=\theta$, i.e.\
$\theta\in Z_{G}(\sigma)$, the centralizer of $\sigma$ in $\Aut(\g)$.  More
generally, allowing the inversion piece, a constant gauge automorphism is
admissible iff $\theta\in N_\Gamma(\g)=\Aut(\g,\sigma):=
\{\theta\in\Aut(\g):\theta\langle\sigma\rangle\theta^{-1}=
\langle\sigma\rangle\}$ (those $\theta$ that, together with a suitable
$\varphi$, satisfy \eqref{eq:equiv}); these are item (1) of
Theorem~\ref{thm:loop}.  This proves the constant part.

\emph{Non-constant loops are inner.}  Decompose $G=\Aut(\g)$ with identity
component $G^\circ=\mathrm{Inn}(\g)$ and finite $\mathrm{Out}(\g)$.  Since
$\Gm$ is connected and $\mathrm{Out}(\g)$ is discrete, any loop
$\Psi\colon\Gm\to G$ lands in a single coset $\theta_0 G^\circ$; writing
$\Psi=\theta_0\cdot\Psi^\circ$ with $\Psi^\circ\colon\Gm\to G^\circ=
\mathrm{Inn}(\g)$ separates an outer ``constant'' part $\theta_0$ (item
(1)) from an inner loop $\Psi^\circ$.  Because $\g$ is simple,
$\mathrm{Inn}(\g)=G^\circ$ is a connected semisimple adjoint group and
$\Psi^\circ$ is a Laurent loop into it; the constraint
$\Psi^\circ(\zeta x)=\sigma\Psi^\circ(x)\sigma^{-1}$ defines a
$\sigma$-equivariant inner loop.  These inner loops are exactly the
images of $L(\g,\sigma)$ acting on itself by the exponentiated inner
derivations $\exp(\mathrm{ad}\,\xi)$, $\xi\in L(\g,\sigma)$ nilpotent
elementwise, generating $\mathrm{Inn}(L(\g,\sigma))$; thus the inner gauge
loops are precisely $\mathrm{Inn}(L(\g,\sigma))$.  Combining,
\[
\mathrm{Mor}_\Gamma(\Gm,N_\Gamma(\g))=
\mathrm{Inn}(L(\g,\sigma))\rtimes \Aut(\g,\sigma)/\!\langle\sigma\rangle,
\]
the last factor being the outer/constant data modulo the redundancy that
$\sigma$ itself acts as an inner gauge transformation of the loop algebra
(it is realized by the deck transformation $\gamma\in\Gamma$, already
absorbed by the variety part / definition of $\M^\Gamma$).

\subsection{Assembling Theorem~\ref{thm:loop}}
Collecting the variety part $k^\times\rtimes\Z/2$ and the gauge part, every
$\Phi\in\Aut(L)$ has the form \eqref{eq:loopaut}
$(\Phi f)(t)=\Psi(t)\,f(at^{\pm1})$.  The semidirect/extension structure is
exactly as displayed in Theorem~\ref{thm:loop}:
\begin{itemize}
\item the rescalings $\varphi_a:t\mapsto at$ ($a\in k^\times$) act by
$f(t)\mapsto f(a^{-1}t)$ and give the $k^\times$;
\item the inversion $\iota:t\mapsto t^{-1}$ gives the $\Z/2$, and forces
$\sigma\mapsto\sigma^{-1}$ on $\g$ (computed above), consistent with the
isomorphism $L(\g,\sigma)\cong L(\g,\sigma^{-1})$;
\item the constant outer/diagram part $\theta\in\Aut(\g,\sigma)$ realizes
the symmetries of $(\g,\sigma)$;
\item the inner loops are $\mathrm{Inn}(L)$.
\end{itemize}
Passing to outer automorphisms (quotient by $\mathrm{Inn}(L)$ and by the
variety group $k^\times\rtimes\Z/2$) leaves exactly the symmetries of the
pair $(\g,\langle\sigma\rangle)$ up to conjugacy, i.e.\
$\mathrm{Stab}(\bar\sigma)\subseteq\mathrm{Out}(\g)$, giving the short exact
sequence in Theorem~\ref{thm:loop}.  \qed

\section{Multiloop algebras: proof of Theorem~\ref{thm:multiloop}}
\label{sec:multiloop}

Here $X=(\Gm)^n=\Spec k[t_1^{\pm1},\dots,t_n^{\pm1}]$,
$\Gamma=\prod_{j=1}^n\Z/m_j=\langle\gamma_1,\dots,\gamma_n\rangle$, with
$\gamma_j\cdot t_i=\zeta_{m_j}^{\delta_{ij}}t_i$ and $\gamma_j\cdot=
\sigma_j$ on $\g$, the $\sigma_j$ commuting of orders $m_j$.  Then
$A=k[t^{\pm1}]$ is reduced and finitely generated (H1); the action of
$\Gamma$ on $(k^\times)^n$ is free because $\prod_j\zeta_{m_j}^{a_j}=1$
componentwise forces $m_j\mid a_j$ (H2); $\g$ is simple (H3).
Theorem~\ref{thm:main} applies.

\subsection{Identification with the standard multiloop algebra}
Decompose $\g=\bigoplus_{\bar j\in\Gamma^\vee}\g_{\bar j}$ into
simultaneous eigenspaces of the commuting $\sigma_1,\dots,\sigma_n$, where
$\Gamma^\vee=\prod_j\Z/m_j$ and $\sigma_j|_{\g_{\bar k}}=
\zeta_{m_j}^{k_j}$.  As in the loop case, a Laurent monomial
$a\,t_1^{i_1}\cdots t_n^{i_n}$ ($a\in\g$) is $\Gamma$-invariant iff
$\sigma_j(a)=\zeta_{m_j}^{i_j}a$ for all $j$, i.e.\
$a\in\g_{(\bar i_1,\dots,\bar i_n)}$.  Hence
\[
\M^\Gamma=M(\g,\sigma_1,\dots,\sigma_n)
=\bigoplus_{i\in\Z^n}\g_{(\bar i_1,\dots,\bar i_n)}\,t_1^{i_1}\cdots
t_n^{i_n},
\]
the standard multiloop algebra.

\subsection{The variety automorphisms $\Aut((\Gm)^n,\Gamma)$}
$\Aut_{\mathrm{var}}((\Gm)^n)=(k^\times)^n\rtimes\GL_n(\Z)$:
a variety automorphism is a monomial map
$t_i\mapsto a_i\prod_\ell t_\ell^{B_{i\ell}}$ with $a\in(k^\times)^n$ and
$B\in\GL_n(\Z)$ (these are exactly the $k$-algebra automorphisms of
$k[t^{\pm1}]$, since the units of $k[t^{\pm1}]$ are $k^\times\cdot
(\text{monomials})$ and an automorphism permutes a generating set of
units up to the monomial group; $B$ invertible over $\Z$ is forced by
invertibility).  The normalization condition: $\varphi$ normalizes the
$\Gamma$-action of rotations by roots of unity.  The rotation subgroup of
$\Aut((\Gm)^n)$ generated by $\Gamma$ is the finite group of
diagonal $\zeta$-scalings; $\varphi$ normalizes it iff the matrix $B$ maps
the $\Gamma$-rotation lattice to itself.  Concretely
$\varphi\gamma_j\varphi^{-1}$ is again a $\Gamma$-rotation iff $B$ induces
an automorphism $c\in\Aut(\Gamma)$; thus
\[
\Aut((\Gm)^n,\Gamma)=
\bigl((k^\times)^n\rtimes\GL_n(\Z)\bigr)_{\Gamma}
:=\{(a,B):B\text{ normalizes the }\Gamma\text{-rotation, with induced }
c_B\in\Aut(\Gamma)\}.
\]

\subsection{The gauge maps and the automorphism group}
By Theorem~\ref{thm:main} and Remark~\ref{rem:groupstr},
\[
\Aut\bigl(M(\g,\sigma_1,\dots,\sigma_n)\bigr)\cong
\Aut((\Gm)^n,\Gamma)\ \ltimes\
\mathrm{Mor}_\Gamma\bigl((\Gm)^n,N_\Gamma(\g)\bigr),
\]
where the gauge group consists of regular maps $\Psi\colon(\Gm)^n\to
N_\Gamma(\g)\subseteq\Aut(\g)$ satisfying, for $c=c_B\in\Aut(\Gamma)$,
\[
\Psi(\gamma x)=c(\gamma)\,\Psi(x)\,\gamma^{-1}\qquad(\gamma\in\Gamma),
\]
i.e.\ $\Psi$ is equivariant for the $\Gamma$-rotation on the source and the
$c$-twisted conjugation by $\langle\sigma_1,\dots,\sigma_n\rangle$ on the
target.  As in \S\ref{sec:loop}, since $(\Gm)^n$ is connected and
$\mathrm{Out}(\g)$ is finite, each $\Psi$ factors as
$\Psi=\theta_0\cdot\Psi^\circ$ with $\theta_0$ a constant element of
$N_\Gamma(\g)$ realizing a symmetry of the family
$\{\sigma_1,\dots,\sigma_n\}$ and $\Psi^\circ\colon(\Gm)^n\to
\mathrm{Inn}(\g)$ an inner loop; the inner loops are exactly
$\mathrm{Inn}(M)$.  We write
\[
\Aut(\g;\sigma_1,\dots,\sigma_n):=N_\Gamma(\g)
=\{\theta\in\Aut(\g):\theta\langle\sigma_1,\dots,\sigma_n\rangle
\theta^{-1}=\langle\sigma_1,\dots,\sigma_n\rangle\},
\]
the group of automorphisms of $\g$ permuting/powering the $\sigma_i$
compatibly with some admissible monomial $\varphi$ (the matrix $B$
realizing the same permutation/powering on $\Gamma^\vee$).  Assembling, every
automorphism of $M$ has the form $(\Phi f)(t)=\Psi(t)\,f(\varphi^{-1}t)$
with $\varphi$ monomial normalizing $\Gamma$ and $\Psi$ as above; the
group is
\[
\Aut(M)\cong
\bigl((k^\times)^n\rtimes\GL_n(\Z)\bigr)_\Gamma\ \ltimes\
\Bigl(\,\mathrm{Inn}(M)\ \rtimes\ \Aut(\g;\sigma_1,\dots,\sigma_n)
/\langle\sigma_1,\dots,\sigma_n\rangle\Bigr),
\]
which is the explicit description claimed in
Theorem~\ref{thm:multiloop}. \qed

\section{Relation to the isomorphism classification}\label{sec:iso}

The same rigidity inputs (Lemmas~\ref{lem:centroid}--\ref{lem:fibrerigid})
classify \emph{isomorphisms} $\M(X,\g)^\Gamma\xrightarrow{\sim}
\M(X',\g')^{\Gamma'}$, not just automorphisms: replacing $\Phi$ by an
isomorphism between two equivariant map algebras, the centroid argument
gives a variety isomorphism $X/\Gamma\cong X'/\Gamma'$, the fibre argument
gives $\g\cong\g'$, and the equivariance \eqref{eq:equiv} matches the
$\Gamma$- and $\Gamma'$-structures.  Hence:

\begin{corollary}\label{cor:iso}
Under \textnormal{(H1)--(H3)} for both sides, two equivariant map algebras
$\M(X,\g)^\Gamma$ and $\M(X',\g')^{\Gamma'}$ are isomorphic as Lie algebras
if and only if there is a variety isomorphism $\varphi\colon X\to X'$
intertwining the $\Gamma$- and $\Gamma'$-actions up to an automorphism
$c\colon\Gamma\xrightarrow{\sim}\Gamma'$, together with a regular map
$\Psi\colon X\to\mathrm{Isom}_{\mathrm{Lie}}(\g',\g)$ satisfying the
twisted equivariance \eqref{eq:equiv}.  In particular:
\begin{itemize}
\item Two twisted loop algebras $L(\g,\sigma)$, $L(\g',\sigma')$ are
isomorphic iff $\g\cong\g'$ and $\langle\sigma\rangle$,
$\langle\sigma'\rangle$ are conjugate in $\Aut(\g)$ up to the order
$m\mapsto m$ data; equivalently iff the diagrams and the diagram-automorphism
classes agree (recovering the standard list of affine Kac--Moody
realizations $X_N^{(r)}$).
\item Two multiloop algebras $M(\g,\sigma_\bullet)$,
$M(\g',\sigma'_\bullet)$ are isomorphic iff $\g\cong\g'$ and the subgroups
$\langle\sigma_1,\dots,\sigma_n\rangle$,
$\langle\sigma'_1,\dots,\sigma'_n\rangle$ of $\Aut(\g)$ are conjugate
\emph{after} an admissible change of coordinates $B\in\GL_n(\Z)$ on the
grading lattice; this is the classification of multiloop algebras by the
$\GL_n(\Z)$-conjugacy class of the finite abelian subgroup
$\langle\sigma_\bullet\rangle\subseteq\Aut(\g)$.
\end{itemize}
\end{corollary}
\begin{proof}
Run the proof of Theorem~\ref{thm:main} with an isomorphism in place of an
automorphism: every step (perfectness, centroid, fibre, equivariance) used
only that the map is a Lie isomorphism, and produces the pair
$(\varphi,\Psi)$ relating the two structures.  Conversely such a pair
induces an isomorphism by the computation in the proof of
Theorem~\ref{thm:main}.  Specializing $X=X'=\Gm$ (resp.\ $(\Gm)^n$) and
reading off the constraints \eqref{eq:equiv} for the rotation actions gives
the conjugacy statements: for $\Gm$, $\varphi=t\mapsto at^{\pm1}$ forces
$\sigma'\sim\sigma^{\pm1}$ and $m=m'$; for $(\Gm)^n$, the matrix
$B\in\GL_n(\Z)$ implements the lattice change of coordinates, so the
grading groups $\langle\sigma_\bullet\rangle$ must be $\GL_n(\Z)$-conjugate
in $\Aut(\g)$.  The loop-algebra statement is the classical correspondence
between $\langle\sigma\rangle$-conjugacy classes (i.e.\ diagram
automorphisms of order $r$) and affine types $X_N^{(r)}$.
\end{proof}

\begin{remark}
The hypotheses (H1)--(H3) cover all the listed special cases.  When $\fa$
is reductive but not simple, or $\Gamma$ does not act freely on $X(k)$, the
same method applies after stratifying $X$ by orbit type and replacing the
single simple quotient in Lemma~\ref{lem:fibres} by a product of simple
quotients indexed by the orbit; the centroid then equals
$\OO(X/\Gamma)\otimes Z(\fa)$ and the gauge group is taken with values in
$\Aut(\fa)$.  The structure
``$\Aut=\,$(variety automorphisms normalizing $\Gamma$)$\,\ltimes\,$(twisted
equivariant gauge group)'' persists; this is the broad-generality answer.
\end{remark}

\end{document}
